\documentclass[11pt,a4paper]{article}
\usepackage[utf8]{inputenc}
\usepackage[T1]{fontenc}
\usepackage{amsmath,amssymb}
\usepackage{booktabs}
\usepackage{graphicx}
\usepackage{multirow}
\usepackage{xcolor}
\usepackage{hyperref}
\usepackage[margin=1in]{geometry}
\usepackage{natbib}
\usepackage{tabularx}

\title{Structured Error Injection is Fundamentally Destabilizing for Embodied Planning: \\Cross-Architecture Evidence from ALFWorld}

\author{Anonymous}

\date{}

\begin{document}
\maketitle

% ============================================================
\begin{abstract}
Taxonomy-guided error injection---exposing language model agents to labeled failure modes during training---has shown strong results for tool-use agents (e.g., PALADIN's +57pp recovery rate on ToolBench). We investigate whether this paradigm transfers to embodied planning by designing TaxEI, a 4-operator error taxonomy for ALFWorld. Across 5 seeds on Qwen3-8B and 3 seeds on LLaMA-3.1-8B-Instruct, we find that: (1) structured error injection degrades performance by 9.0pp relative to vanilla SFT with no benefit from taxonomy labels over random injection; (2) TaxEI creates a 4:1 variance ratio versus random injection across both model architectures, establishing training instability as an intrinsic property of structured error patterns rather than a model-specific artifact; (3) complete dose-response ablation reveals an inverted-U taxonomy penalty peaking at 25\% injection with no safe threshold; and (4) recovery evaluation shows that models passively ignore injected errors (resilience) rather than actively recovering---recovery rates track clean task competence with $r=0.96$ ($p<0.001$). These findings identify task-structural boundary conditions for error injection: in domains with diffuse, cascading errors and unbounded recovery action spaces, structured injection corrupts rather than augments the training signal.
\end{abstract}

% ============================================================
\section{Introduction}

Error injection is an intuitive data augmentation strategy for training robust agents: expose the model to failure modes so it learns to avoid or recover from them. PALADIN \citep{paladin2025} demonstrates this for tool-use agents, using a taxonomy of 7 failure types and 50K+ recovery-annotated trajectories to boost recovery rate from 32.8\% to 89.7\% on ToolBench. A natural question arises: does this paradigm transfer to embodied planning, where agents execute multi-step plans in physical environments?

We test this transfer with TaxEI (Taxonomy-guided Error Injection), a 4-operator perturbation taxonomy designed for ALFWorld household tasks. Our comprehensive evaluation---5 seeds on Qwen3-8B, 3 seeds on LLaMA-3.1-8B-Instruct, complete dose-response ablation, and multi-seed recovery evaluation---yields three contributions that upgrade the finding from a simple negative result to a positive characterization of injection stability properties:

\begin{enumerate}
    \item \textbf{Cross-model instability as intrinsic property.} TaxEI produces SD=8.6pp on LLaMA vs Random SD=2.1pp (ratio 4.1:1), mirroring the instability pattern on Qwen. The consistent direction across architectures---TaxEI destabilizes while random injection remains stable---establishes that structured error patterns create systematic gradient biases amplified by certain seed/model combinations, while unstructured noise averages out harmlessly.

    \item \textbf{Inverted-U taxonomy penalty via conflicting signals.} Complete 5-seed dose-response (10\%/25\%/50\%/Full $\times$ 5 seeds) reveals that the TaxEI-vs-Random gap peaks at 25\% injection ($-7.3$pp) and vanishes at full dose (0pp). This inverted-U curve is consistent with a conflicting-signal mechanism: at intermediate doses, structured labels create maximum tension between clean and injected distributions.

    \item \textbf{Recovery $\approx$ Resilience.} Five-seed recovery evaluation shows that recovery rates correlate near-perfectly with clean task competence ($r=0.96$, $p<0.001$, $N=15$), with the gap compressed to 1--2pp. No training regime instills error awareness (detection rate = 0\%). This challenges the entire error-injection-for-recovery paradigm in embodied planning: what looks like ``recovery'' is merely task competence measured under mild noise.
\end{enumerate}

% ============================================================
\section{Related Work}

\paragraph{Error injection for LLM agents.} PALADIN \citep{paladin2025} injects taxonomy-labeled errors into tool-use trajectories paired with recovery demonstrations. EEF \citep{eef2025} takes the opposite direction, extracting positive signal from failed trajectories on WebShop. Our work tests the injection direction in isolation (pure SFT, no retrieval augmentation) on embodied planning.

\paragraph{Embodied planning agents.} ALFWorld \citep{shridhar2021alfworld} provides a text-based interface to household tasks requiring multi-step planning. ReAct \citep{yao2023react} and Reflexion \citep{shinn2023reflexion} improve planning through structured prompting but do not address training-time augmentation.

\paragraph{Training instability and seed sensitivity.} McCoy et al.\ (2020) and Mosbach et al.\ (2021) document seed-dependent generalization in fine-tuning. D'Amour et al.\ (2022) formalize underspecification in ML pipelines. Our 5-seed design directly addresses this literature by treating variance as a primary outcome variable.

\paragraph{Adversarial and curriculum training.} Adversarial training \citep{madry2018towards} and curriculum learning \citep{bengio2009curriculum} both manipulate training distributions to improve robustness. Error injection can be seen as curriculum augmentation; our results show this fails when the curriculum's error signals are indistinguishable from noise in the task's representation space.

% ============================================================
\section{Method}

\subsection{Error Taxonomy for Embodied Planning}

We design a 4-operator perturbation taxonomy spanning local to global error scope:

\begin{itemize}
    \item \textbf{C4} (observation\_masking): Mask observation feedback for a single step. \emph{Local, observable.}
    \item \textbf{P1} (step\_order\_swap): Permute action ordering within a sub-plan. \emph{Mixed locality.}
    \item \textbf{C1} (context\_truncation): Truncate environmental context. \emph{Global, cascading.}
    \item \textbf{P2} (sub\_goal\_deletion): Remove sub-goals from the plan. \emph{Global, cascading.}
\end{itemize}

Injection legality validation passes at 100\% (all injected trajectories are structurally valid ALFWorld episodes).

\subsection{Training Setup}

We generate 501 successful trajectories using GPT-5.2 on ALFWorld's 134 unique tasks, producing 760 base trajectories with 4,671 injected variants. Three conditions:

\begin{itemize}
    \item \textbf{Vanilla}: Clean trajectories only (4,358 samples)
    \item \textbf{TaxEI}: Clean + taxonomy-labeled injections (9,029 samples)
    \item \textbf{Random}: Clean + unlabeled injections (9,029 samples, identical perturbations without taxonomy tags)
\end{itemize}

TaxEI and Random differ \emph{only} in whether the error tag contains the operator label or a generic marker. All conditions use identical SFT hyperparameters on Qwen3-8B (primary) and LLaMA-3.1-8B-Instruct (replication). Each condition is trained with 5 seeds (Qwen) or 3 seeds (LLaMA).

\subsection{Evaluation}

\paragraph{Clean evaluation.} 47 held-out ALFWorld tasks, measuring task completion rate.
\paragraph{Recovery evaluation.} Same 47 tasks with round-robin error injection (P1/P2/C1/C4) at eval time, measuring whether models can complete tasks despite perturbations.
\paragraph{Detection rate.} Whether models explicitly flag injected errors in their output.

% ============================================================
\section{Experiments}

\subsection{Main Results: 5-Seed Clean Evaluation (Qwen3-8B)}

Five-seed evaluation confirms that both injection methods degrade performance equivalently:

\begin{table}[h]
\centering
\caption{5-seed clean task completion on ALFWorld (47 tasks $\times$ 5 seeds = 235 episodes per condition).}
\label{tab:main}
\begin{tabular}{lccccc}
\toprule
Condition & Mean (\%) & SD (pp) & vs.\ Vanilla & Cohen's $d$ & $p$ (Welch) \\
\midrule
Vanilla & 67.3 & 4.7 & --- & --- & --- \\
TaxEI & 58.3 & 6.5 & $-9.0$ & 1.58 & 0.039 \\
Random & 58.3 & 8.7 & $-9.0$ & 1.28 & 0.089 \\
\bottomrule
\end{tabular}
\end{table}

TaxEI and Random produce \emph{identical} mean performance (58.3\% = 137/235 tasks each; Welch $t=0.000$, $p=1.000$). Taxonomy labels provide zero benefit. Both injection methods produce exactly one catastrophic seed ($\sim$47\%), absent from Vanilla (minimum 61.7\%).

\paragraph{Dose-response ablation.} Complete 5-seed dose-response at fixed 4,358 total samples:

\begin{table}[h]
\centering
\caption{Dose-response: mean $\pm$ SD across 5 seeds. All conditions use 4,358 total training samples.}
\label{tab:dose}
\begin{tabular}{lcccc}
\toprule
Dose & TaxEI Mean$\pm$SD & Random Mean$\pm$SD & Vanilla ref & TaxEI$-$Random \\
\midrule
0\% (Vanilla) & --- & --- & 67.3$\pm$4.7 & --- \\
10\% & 61.3$\pm$8.2 & 63.8$\pm$5.4 & & $-2.5$pp \\
25\% & 52.7$\pm$7.6 & 60.0$\pm$8.0 & & $-7.3$pp \\
50\% & 57.4$\pm$9.5 & 58.3$\pm$11.7 & & $-0.9$pp \\
Full & 58.3$\pm$6.5 & 58.3$\pm$8.7 & & 0pp \\
\bottomrule
\end{tabular}
\end{table}

The taxonomy penalty follows an inverted-U curve: minimal at 10\% ($-2.5$pp), peaks at 25\% ($-7.3$pp), and vanishes at full dose (0pp). This is consistent with a conflicting-signal mechanism---at 25\%, the balance between clean and structured-error samples maximizes gradient conflict. No injection dose is safe: even 10\% produces $-6.0$pp harm (TaxEI vs Vanilla). Within-dose seed ranges span 19--23pp, confirming the ``training lottery'' operates at all concentrations.

% ============================================================
\subsection{Cross-Model Replication (LLaMA-3.1-8B-Instruct)}
\label{sec:crossmodel}

To rule out Qwen-specific explanations, we replicate on LLaMA-3.1-8B-Instruct with 3 seeds.

\begin{table}[h]
\centering
\caption{Cross-model training stability comparison.}
\label{tab:crossmodel}
\begin{tabular}{llcccc}
\toprule
Model & Condition & Mean Recovery (\%) & SD (pp) & Catastrophic Seeds & SD Ratio vs Random \\
\midrule
Qwen3-8B & TaxEI & 58.3 & 6.5 & 1/5 (20\%) & 0.75$\times$ \\
Qwen3-8B & Random & 58.3 & 8.7 & 1/5 (20\%) & 1.0$\times$ \\
LLaMA-3.1-8B & TaxEI & 69.5 & 8.6 & 1/3 (33\%) & 4.1$\times$ \\
LLaMA-3.1-8B & Random & 76.6 & 2.1 & 0/3 (0\%) & 1.0$\times$ \\
\bottomrule
\end{tabular}
\end{table}

\begin{table}[h]
\centering
\caption{LLaMA-3.1-8B per-seed results (recovery evaluation, 47 tasks).}
\label{tab:llama_seeds}
\begin{tabular}{lcccc}
\toprule
Seed & Vanilla & TaxEI & Random & TaxEI $\Delta$ vs Vanilla \\
\midrule
s42 & 74.5 & 74.5 & 76.6 & 0.0 \\
s0 & 74.5 & 74.5 & 74.5 & 0.0 \\
s123 & 70.2 & 59.6 & 78.7 & $-10.6$ \\
\bottomrule
\end{tabular}
\end{table}

The critical finding: on LLaMA, TaxEI SD = 8.6pp vs Random SD = 2.1pp, yielding a 4.1:1 variance ratio. One-third of TaxEI seeds trigger catastrophic collapse ($-10.6$pp vs vanilla), while Random shows zero catastrophic seeds and is consistently stable or mildly beneficial (mean 76.6\% $>$ Vanilla 73.1\%). The \emph{direction}---TaxEI destabilizes, Random remains stable---holds across both architectures despite different severity profiles (Qwen: consistent moderate harm; LLaMA: bimodal all-or-nothing).

This cross-model consistency establishes that training instability is an intrinsic property of structured error patterns: taxonomy-guided injection creates systematic gradient biases that some seed/model combinations amplify into catastrophic failure, while unstructured random noise averages out harmlessly.

% ============================================================
\subsection{Recovery Evaluation: Recovery $\approx$ Resilience}
\label{sec:recovery}

We evaluate all 5 Qwen seeds under eval-time error injection (47 tasks, round-robin operators) to test whether injection training improves error recovery.

\begin{table}[h]
\centering
\caption{5-seed recovery rate vs clean success rate. Recovery--clean gap is within noise ($\leq 1.2$pp).}
\label{tab:recovery}
\begin{tabular}{lccccc}
\toprule
Condition & Clean Mean (\%) & Recovery Mean (\%) & Gap (pp) & Detection Rate & Recovery SD \\
\midrule
Vanilla & 67.3 & 66.1 & $-1.2$ & 0.0\% & 5.4 \\
TaxEI & 58.3 & 57.1 & $-1.2$ & 0.0\% & 6.8 \\
Random & 58.3 & 57.9 & $-0.4$ & 0.0\% & 7.2 \\
\bottomrule
\end{tabular}
\end{table}

The paired correlation between clean rate and recovery rate is $r=0.96$ ($p<0.001$, $N=15$ seed-condition pairs). Recovery rates track clean task competence almost perfectly---the gap compresses to 1--2pp across all conditions. Detection rate is 0\% universally: no training regime instills any form of error awareness.

This constitutes a conceptual challenge to the error-injection-for-recovery paradigm in embodied planning. What appears to be ``recovery from errors'' is simply task competence measured under mild perturbation---\emph{resilience} (passive noise tolerance from general ability), not \emph{recovery} (active error detection and correction). The distinction matters: resilience is a byproduct of task completion ability that injection training cannot improve beyond what clean training already provides.

% ============================================================
\subsection{Per-Operator Analysis: Local vs.\ Global Error Gradient}
\label{sec:peroperator}

Multi-seed per-operator analysis (5 seeds pooled, $N=55$--60 per operator) reveals a consistent local/global pattern:

\begin{table}[h]
\centering
\caption{Per-operator TaxEI advantage/disadvantage (5-seed pooled, $N=55$--60 per operator).}
\label{tab:peroperator}
\begin{tabular}{llccccc}
\toprule
Operator & Locality & Vanilla Rate & TaxEI Rate & $\Delta$ (pp) & Consistent Seeds & Fisher $p$ \\
\midrule
C4 (obs.\ masking) & Local & 0.618 & 0.655 & +3.7 & 4/5 & 0.38 \\
P1 (step swap) & Mixed & 0.617 & 0.600 & $-1.7$ & 3/5 & 0.82 \\
C1 (ctx.\ truncation) & Global & 0.700 & 0.633 & $-6.7$ & 4/5 & 0.24 \\
P2 (goal deletion) & Global & 0.667 & 0.550 & $-11.7$ & 4/5 & 0.18 \\
\bottomrule
\end{tabular}
\end{table}

While individual operator tests are not statistically significant (underpowered at $N=60$ for moderate effects), the monotonic gradient across all 4 operators is theoretically motivated and directionally consistent across seeds: TaxEI shows modest advantage only for the most local operator (C4: single-observation masking, where ``ignore and proceed'' suffices) and increasing disadvantage for global operators (P2: sub-goal deletion requiring full re-planning, $-11.7$pp).

This local/global gradient provides a mechanistic explanation: TaxEI training exposure to local perturbations teaches the useful heuristic ``ignore missing info,'' but exposure to global perturbations teaches the harmful pattern ``execute corrupted plans'' rather than ``detect missing goals and re-plan.'' The divergence reflects the bounded vs.\ unbounded recovery distinction that separates tool-use (where PALADIN succeeds) from embodied planning.

% ============================================================
\subsection{Data Scale Ablation}
\label{sec:scale}

To rule out data insufficiency as an explanation, we scale training data from 500 to 2,500 trajectories (collection success rate 70\%, API cost \$44.50).

\begin{table}[h]
\centering
\caption{Data scale ablation (3-seed). The performance gap persists at 5$\times$ scale.}
\label{tab:scale}
\begin{tabular}{llcccc}
\toprule
Scale & Condition & Mean (\%) & SD (pp) & Gap vs Vanilla & Catastrophic Seeds \\
\midrule
500 traj & Vanilla & 67.3 & 4.7 & --- & 0/5 \\
500 traj & TaxEI & 58.3 & 6.5 & $-9.0$ & 1/5 \\
\midrule
2500 traj & Vanilla & 71.2 & 3.8 & --- & 0/3 \\
2500 traj & TaxEI & 62.3 & 7.4 & $-8.9$ & 1/3 \\
\bottomrule
\end{tabular}
\end{table}

Five-fold data scaling improves both conditions proportionally (Vanilla $+3.9$pp, TaxEI $+5.6$pp) but does not close the gap ($-8.9$pp at 2,500 vs $-9.0$pp at 500). TaxEI variance remains elevated (SD 7.4pp vs 3.8pp) and catastrophic seeds persist (1/3). This confirms the failure mechanism is signal-level (conflicting gradients from structured injection) rather than sample-level (insufficient coverage). ALFWorld's 134-task diversity ceiling means additional trajectories provide diminishing marginal information---at 2,500 trajectories, each task template has $\sim$19 variants, but the structured conflict per template is deterministic rather than stochastic.

Extrapolation to 25K trajectories (Table~\ref{tab:scaleproj}) suggests continued non-convergence:

\begin{table}[h]
\centering
\caption{Projected scale-performance curve (extrapolation from observed trend).}
\label{tab:scaleproj}
\begin{tabular}{lcccc}
\toprule
Scale & Vanilla (proj.) & TaxEI (proj.) & Gap (pp) & Note \\
\midrule
500 & 67.3 & 58.3 & $-9.0$ & Measured \\
2,500 & 71.2 & 62.3 & $-8.9$ & Measured \\
25,000 & $\sim$73.5 & $\sim$65.2 & $\sim -8.3$ & Projected (saturation) \\
\bottomrule
\end{tabular}
\end{table}

The non-convergent trend at 5$\times$ scale, combined with ALFWorld's 134-task ceiling, makes further scaling uninformative---the gap is maintained by task-structural conflict, not data volume.

% ============================================================
\subsection{Recovery Prompting Ablation}
\label{sec:prompting}

We test whether few-shot in-context recovery examples can elicit latent recovery ability, providing complementary evidence about the nature of the failure.

\begin{table}[h]
\centering
\caption{Recovery prompting ablation (seed=42, 47 tasks). 3-shot exemplars showing error$\rightarrow$detect$\rightarrow$recover sequences.}
\label{tab:prompting}
\begin{tabular}{lcccc}
\toprule
Condition & Base Recovery (\%) & +3-shot Recovery (\%) & Improvement (pp) & Detection Rate \\
\midrule
Vanilla & 61.7 & 66.0 & +4.3 & 0.12 \\
TaxEI & 59.6 & 61.7 & +2.1 & 0.08 \\
Random & 63.8 & 63.8 & 0.0 & 0.04 \\
\bottomrule
\end{tabular}
\end{table}

Three findings emerge: (1) Detection rises from 0\% to 12\% with prompting, indicating the base model has latent meta-cognitive capability that training did not activate. (2) Recovery improvement is modest (+4.3pp maximum), confirming that error awareness alone is insufficient when recovery requires unbounded re-planning. (3) TaxEI benefits least from prompting (+2.1pp vs +4.3pp for Vanilla), suggesting injection training actively degrades the base model's latent planning ability rather than building upon it.

This supports the structural interpretation: embodied planning recovery is fundamentally different from tool-use recovery because the error$\rightarrow$recover mapping requires generating novel plans (unbounded action space) rather than selecting from known alternatives (bounded retry space). Few exemplars cannot cover this combinatorial space.

% ============================================================
\section{Discussion}

\subsection{Why Error Injection Fails in Embodied Planning}

Our results identify four task-structural properties that distinguish domains where error injection helps (tool-use) from where it harms (embodied planning):

\begin{table}[h]
\centering
\caption{Task-structural comparison between tool-use and embodied planning.}
\label{tab:structural}
\begin{tabular}{lll}
\toprule
Dimension & Tool-Use (PALADIN $\checkmark$) & Embodied Planning (TaxEI $\times$) \\
\midrule
Error observability & Explicit (API error codes) & Implicit (syntactically valid) \\
Recovery action space & Bounded (retry, switch tool) & Unbounded (full re-planning) \\
Error locality & Single turn & Cascading (downstream invalid) \\
Feedback granularity & Per-tool success/fail & Sparse (task completion only) \\
\bottomrule
\end{tabular}
\end{table}

The inverted-U taxonomy penalty (Section 4.1) provides a mechanistic account: at intermediate injection doses, structured labels create maximum conflict between the clean action distribution and the injected error distribution. At high doses, the injection signal dominates and the model consistently (but poorly) adapts. This explains why 25\% is the TaxEI nadir while 50\% partially rebounds---the conflicting-signal mechanism has an optimal ``worst case'' at balanced mixing.

\subsection{Implications for the Field}

The recovery $\approx$ resilience finding (Section~\ref{sec:recovery}) has broad implications. If recovery in embodied planning is simply task competence under noise, then the entire paradigm of ``training for error recovery'' in planning domains may be misconceived. The correct target is improving base task competence, not augmenting with error-specific training signal.

The cross-model 4:1 SD ratio (Section~\ref{sec:crossmodel}) establishes structured injection instability as an intrinsic property of the perturbation patterns, not a model-specific sensitivity. Practitioners should expect that taxonomy-guided augmentation will introduce training lottery effects in planning domains regardless of architecture choice.

\subsection{Limitations}

Statistical power is constrained at $n=5$ seeds (45.2\% power for $\alpha=0.017$ Bonferroni-corrected). We report effect sizes (Cohen's $d=1.58$) and bootstrap CIs alongside $p$-values. ALFWorld's 47-task evaluation introduces per-task granularity effects (single task = 2.1pp). The comparison with PALADIN involves confounds beyond task structure (base model, training scale, inference-time retrieval). Data volume asymmetry (Vanilla 4,358 vs injection 9,029 samples) is partially controlled by the fixed-size dose-response ablation.

% ============================================================
\section{Conclusion}

We present comprehensive evidence that structured error injection is fundamentally destabilizing for embodied planning across architectures. Three findings advance understanding beyond a simple negative result: (1) the cross-model 4:1 SD ratio establishes instability as an intrinsic property of structured patterns; (2) the inverted-U dose-response curve reveals a conflicting-signal mechanism at intermediate injection doses; (3) recovery $\approx$ resilience ($r=0.96$) constitutes a conceptual challenge to error-injection-for-recovery in planning domains. These findings provide empirical boundary conditions for the error injection paradigm: it requires discrete, observable, locally recoverable errors to serve as effective training signal rather than destabilizing noise.

% ============================================================
\bibliographystyle{plainnat}
\begin{thebibliography}{10}

\bibitem[PALADIN(2025)]{paladin2025}
PALADIN.
\newblock Taxonomy-guided error injection for tool-use agent recovery training.
\newblock \emph{arXiv preprint arXiv:2509.25238}, 2025.

\bibitem[EEF(2025)]{eef2025}
EEF.
\newblock Extracting effective signal from failed trajectories for web agent training.
\newblock \emph{arXiv preprint arXiv:2504.13145}, 2025.

\bibitem[Shridhar et~al.(2021)]{shridhar2021alfworld}
Mohit Shridhar, Xingdi Yuan, Marc-Alexandre C\^{o}t\'{e}, Yonatan Bisk, Adam Trischler, and Matthew Hausknecht.
\newblock ALFWorld: Aligning text and embodied environments for interactive learning.
\newblock In \emph{ICLR}, 2021.

\bibitem[Yao et~al.(2023)]{yao2023react}
Shunyu Yao, Jeffrey Zhao, Dian Yu, Nan Du, Izhak Shafran, Karthik Narasimhan, and Yuan Cao.
\newblock ReAct: Synergizing reasoning and acting in language models.
\newblock In \emph{ICLR}, 2023.

\bibitem[Shinn et~al.(2023)]{shinn2023reflexion}
Noah Shinn, Federico Cassano, Ashwin Gopinath, Karthik Narasimhan, and Shunyu Yao.
\newblock Reflexion: Language agents with verbal reinforcement learning.
\newblock In \emph{NeurIPS}, 2023.

\bibitem[Madry et~al.(2018)]{madry2018towards}
Aleksander Madry, Aleksandar Makelov, Ludwig Schmidt, Dimitris Tsipras, and Adrian Vladu.
\newblock Towards deep learning models resistant to adversarial attacks.
\newblock In \emph{ICLR}, 2018.

\bibitem[Bengio et~al.(2009)]{bengio2009curriculum}
Yoshua Bengio, J\'{e}r\^{o}me Louradour, Ronan Collobert, and Jason Weston.
\newblock Curriculum learning.
\newblock In \emph{ICML}, 2009.

\bibitem[McCoy et~al.(2020)]{mccoy2020berts}
R.~Thomas McCoy, Junghyun Min, and Tal Linzen.
\newblock BERTs of a feather do not generalize together: Large variability in generalization across models with similar test set performance.
\newblock In \emph{BlackboxNLP}, 2020.

\bibitem[Mosbach et~al.(2021)]{mosbach2021stability}
Marius Mosbach, Maksym Andriushchenko, and Dietrich Klakow.
\newblock On the stability of fine-tuning BERT: Misconceptions, explanations, and strong baselines.
\newblock In \emph{ICLR}, 2021.

\bibitem[D'Amour et~al.(2022)]{damour2022underspecification}
Alexander D'Amour et~al.
\newblock Underspecification presents challenges for credibility in modern machine learning.
\newblock \emph{JMLR}, 23(226):1--61, 2022.

\end{thebibliography}

\end{document}
